% -----------------------------------------------------------------------------
% File: sections/05_conclusion.tex
% Description: Section 5 - Overall Conclusion
% -----------------------------------------------------------------------------

\section{Overall Conclusion}
\label{sec:conclusion}

This report has presented a comprehensive analysis and synthesis of three state-of-the-art research papers presented at ICCCI 2025, addressing key challenges in computer-aided diagnostics and automated software engineering.

\subsection{Summary of Key Findings}

\begin{enumerate}
    \item \textbf{Leukemia Detection (YOLOv11 + GCLI):} The GCLI module effectively replaces the C2PSA self-attention block, achieving the highest mAP50-95 of 25.1\% while reducing over 11\% of parameters. The module's strength lies in its ability to capture global--local feature interactions without dimensionality reduction, making it a lightweight yet powerful attention mechanism for cell-level object detection.

    \item \textbf{Alzheimer Diagnosis (VGG16 Fusion):} The dual-model disjunction fusion strategy achieves 89.87\% accuracy and 87.50\% sensitivity, recovering 5 previously missed AD patients. The paper makes an important methodological contribution by enforcing subject-wise data splitting to eliminate the data leakage problem prevalent in prior studies.

    \item \textbf{LLM Code Review:} Claude 3.5 Sonnet demonstrates the highest alignment with human expert consensus (9/12 metrics), while GPT-4o-mini provides the best cost-to-accuracy ratio at 31.6$\times$ lower cost. The study establishes that LLMs can serve as effective first-pass code review tools, though they cannot yet fully replace senior human reviewers.
\end{enumerate}

\subsection{Cross-Domain Observations}

Comparing the three research domains reveals important parallels and contrasts:
\begin{itemize}
    \item \textbf{Data Integrity:} Both medical imaging and code review face data integrity challenges, though manifested differently. In medical imaging, data leakage (sharing patient slices between train/test splits) artificially inflates performance metrics. In code review, the challenge is context limitation: single-file LLM analysis misses project-level dependencies, yielding incomplete assessments.

    \item \textbf{Sensitivity vs. Precision Trade-off:} In clinical CAD systems, sensitivity is the primary safety metric because false negatives delay critical treatment. The VGG16 fusion paper addresses this through logical disjunction ($L_1 \lor L_2$). Conversely, in software engineering, developers are sensitive to ``linter fatigue'' from excessive false positives. Review systems must balance precision with API costs, making GPT-4o-mini a viable production option.

    \item \textbf{Lightweight vs. Heavyweight Approaches:} All three papers demonstrate that domain-specific, targeted solutions can outperform heavyweight general-purpose methods. GCLI replaces expensive self-attention with lightweight 1D convolutions. Progressive transfer learning adapts pre-trained models without full retraining. And carefully prompted LLMs can approximate human-level code assessment at a fraction of the cost.
\end{itemize}

\subsection{Limitations and Future Directions}
Despite the promising results, each paper presents opportunities for further development:
\begin{itemize}
    \item The GCLI module was evaluated only on YOLOv11; testing on other YOLO variants and two-stage detectors would strengthen generalizability claims.
    \item The VGG16 fusion approach uses 2D slice processing on inherently 3D MRI volumes; 3D architectures could capture volumetric spatial information more effectively.
    \item The LLM code review study evaluates single-file contexts; extending to project-level analysis with dependency-aware prompting could significantly improve SOLID principle assessments.
\end{itemize}

These findings collectively underscore that while deep learning models and LLMs provide strong performance baselines, achieving robust real-world deployment requires careful attention to data integrity, domain-specific architectural design, and rigorous statistical validation.
